
\documentclass{article}
\usepackage{colortbl}
\usepackage{makecell}
\usepackage{multirow}
\usepackage{supertabular}

\begin{document}

\newcounter{utterance}

\centering \large Interaction Transcript for game `wordle', experiment `medium\_frequency\_words\_no\_clue\_no\_critic', episode 4 with reflexion{-}Qwen3.5{-}4B.
\vspace{24pt}

{ \footnotesize  \setcounter{utterance}{1}
\setlength{\tabcolsep}{0pt}
\begin{supertabular}{c@{$\;$}|p{.15\linewidth}@{}p{.15\linewidth}p{.15\linewidth}p{.15\linewidth}p{.15\linewidth}p{.15\linewidth}}
    \# & \multicolumn{2}{c}{Player} && \multicolumn{2}{c}{Game Master} \\
    \hline

    \theutterance \stepcounter{utterance}  
    & & & \multicolumn{4}{p{0.6\linewidth}}{
        \cellcolor[rgb]{0.9,0.9,0.9}{
            \makecell[{{p{\linewidth}}}]{
                \texttt{\tiny{[P1$\langle$GM]}}
                \texttt{You have attempted this same task before and did not succeed. These are the plans you wrote after those attempts.} \\
\texttt{Trial \#1: explanation: The feedback confirms 'a' is first and 'l' is second. 'a', 'c', 'k' are not in the word. I need a 5{-}letter word starting with "al\_" that does not contain 'o', 'u', 'd', 'e', 'p', 'h', 'a', 'c', 'k'. I will try "alms" (4 letters) {-}> "almsy" (no). Let's try "alacr" (no). How about "alate" (t is red). Let's try "alack" (failed). Let's try "aloud" (failed). Let's try "alms" (4 letters). Wait, I need 5 letters. Let's try "alms" + something. "almsy" is not a word. How about "alms" + "e"? "alme" (no). Let's try "alms" + "t"? "almt" (no). Let's try "alms" + "s"? "alms" is 4. Let's try "alms" + "y"? "almsy" (no). Let's try "alms" + "i"? "almsi" (no). Let's try "alms" + "n"? "almsn" (no). Let's try "alms" + "r"? "almsr" (no). Let's try "alms" + "g"? "almsg" (no). Let's try "alms" + "b"? "almsb" (no). Let's try "alms" + "f"? "almsf" (no). Let's try "alms" + "j"? "almsj" (no). Let's try "alms" + "q"? "almsq" (no). Let's try "alms" + "w"? "almsw" (no). Let's try "alms" + "x"? "almsx" (no). Let's try "alms" + "z"? "almsz" (no). Let's try "alms" + "v"? "almsv" (no). Let's try "alms" + "t"? "alms" is 4. Let's try "alms" + "s"? "alms" is 4. Let's try "alms" + "r"? "almsr" (no). Let's try "alms" + "g"? "almsg" (no). Let's try "alms" + "b"? "almsb" (no). Let's try "alms" + "f"? "almsf" (no). Let's try "alms" + "} \\
\\ 
\texttt{Trial \#2: New plan: The previous attempt failed because I generated an explanation that was excessively long, repetitive, and contained hallucinations (e.g., claiming "alace" was a common word when it is not, or claiming "alder" was the answer when the feedback contradicted it). The Game Master rejected the message, likely due to the violation of the "concise" rule or the inclusion of invalid reasoning. I must strictly adhere to the format: a single, clear "explanation:" line followed immediately by a "guess:" line. I will stop guessing "alder" (which failed) and instead use the confirmed pattern "a\_l\_e\_r" (where 'd' is red) to find a valid 5{-}letter word like "alger" (if valid) or switch strategy to a known word fitting the pattern if none exist, ensuring the explanation is brief and the guess is a valid English word. I will also ensure I do not repeat the same word or make up facts about word validity in the explanation.} \\
\\ 
\\ 
\texttt{You are a language wizard who likes to guess words by using the given rules.} \\
\\ 
\texttt{Welcome to Wordle! You have six attempts to guess the target word, a valid English word of five lowercase letters (a{-}z). Please use the tags "explanation:" and "guess:" to provide a concise explanation for each guess.} \\
\\ 
\texttt{For instance, if your guess is "apple", your response should be} \\
\texttt{explanation: this is a common five{-}letter English word, and I am starting my guess with this word.} \\
\texttt{guess: apple} \\
\\ 
\texttt{After each guess, your answer will be validated, and you will receive feedback indicating which letters are correct (green), which letters are correct but in the wrong position (yellow), and which letters are incorrect (red). This feedback can be useful in determining which letters to include or exclude in your next guess.} \\
\\ 
\texttt{For example, the feedback for "apple" might be:} \\
\texttt{guess\_feedback: a<yellow> p<yellow> p<green> l<yellow> e<red>} \\
\\ 
\texttt{The explanation should contain details about how the guess\_feedback is used to arrive at a new guess.} \\
\\ 
\texttt{Do not generate any other text.} \\
\\ 
\texttt{Let's begin with your first guess.} \\
            }
        }
    }
     \\ \\

    \theutterance \stepcounter{utterance}  
    & \multicolumn{4}{p{0.6\linewidth}}{
        \cellcolor[rgb]{0.561,0.737,0.561}{
            \makecell[{{p{\linewidth}}}]{
                \texttt{\tiny{[P1$\rangle$GM]}}
                \texttt{explanation: I will start with a common five{-}letter word to gather initial information about the letters and positions.} \\
\texttt{guess: apple} \\
            }
        }
    }
    & & \\ \\

    \theutterance \stepcounter{utterance}  
    & & & \multicolumn{2}{p{0.3\linewidth}}{
        \cellcolor[rgb]{0.9,0.9,0.9}{
            \makecell[{{p{\linewidth}}}]{
                \texttt{\tiny{[GM$|$GM]}}
                \texttt{attempts = 1} \\
\texttt{target\_word = lager} \\
\texttt{guess = apple} \\
\texttt{guess\_feedback = a<yellow> p<red> p<red> l<yellow> e<yellow>} \\
            }
        }
    }
    & & \\ \\

    \theutterance \stepcounter{utterance}  
    & & & \multicolumn{4}{p{0.6\linewidth}}{
        \cellcolor[rgb]{0.9,0.9,0.9}{
            \makecell[{{p{\linewidth}}}]{
                \texttt{\tiny{[P1$\langle$GM]}}
                \texttt{guess\_feedback: a<yellow> p<red> p<red> l<yellow> e<yellow>} \\
\\ 
\texttt{Provide your response only in this format.} \\
\texttt{explanation: details} \\
\texttt{guess: word} \\
            }
        }
    }
     \\ \\

    \theutterance \stepcounter{utterance}  
    & \multicolumn{4}{p{0.6\linewidth}}{
        \cellcolor[rgb]{0.561,0.737,0.561}{
            \makecell[{{p{\linewidth}}}]{
                \texttt{\tiny{[P1$\rangle$GM]}}
                \texttt{explanation: The feedback 'a<yellow> p<red> p<red> l<yellow> e<yellow>' tells me the word contains 'a', 'l', and 'e', but not 'p'. 'a' is first, 'l' is fourth, and 'e' is fifth. I need a 5{-}letter word starting with "ale\_" where the second letter is not 'p'.} \\
\texttt{guess: alert} \\
            }
        }
    }
    & & \\ \\

    \theutterance \stepcounter{utterance}  
    & & & \multicolumn{2}{p{0.3\linewidth}}{
        \cellcolor[rgb]{0.9,0.9,0.9}{
            \makecell[{{p{\linewidth}}}]{
                \texttt{\tiny{[GM$|$GM]}}
                \texttt{attempts = 2} \\
\texttt{target\_word = lager} \\
\texttt{guess = alert} \\
\texttt{guess\_feedback = a<yellow> l<yellow> e<yellow> r<yellow> t<red>} \\
            }
        }
    }
    & & \\ \\

    \theutterance \stepcounter{utterance}  
    & & & \multicolumn{4}{p{0.6\linewidth}}{
        \cellcolor[rgb]{0.9,0.9,0.9}{
            \makecell[{{p{\linewidth}}}]{
                \texttt{\tiny{[P1$\langle$GM]}}
                \texttt{guess\_feedback: a<yellow> l<yellow> e<yellow> r<yellow> t<red>} \\
\\ 
\texttt{Provide your response only in this format.} \\
\texttt{explanation: details} \\
\texttt{guess: word} \\
            }
        }
    }
     \\ \\

    \theutterance \stepcounter{utterance}  
    & \multicolumn{4}{p{0.6\linewidth}}{
        \cellcolor[rgb]{0.561,0.737,0.561}{
            \makecell[{{p{\linewidth}}}]{
                \texttt{\tiny{[P1$\rangle$GM]}}
                \texttt{explanation: The feedback 'a<yellow> l<yellow> e<yellow> r<yellow> t<red>' confirms the word is "alere" with 't' being incorrect. Since 'a', 'l', 'e', and 'r' are present but 't' is not, and the pattern suggests a common word, I will try "alert" again to confirm or switch to "alder" if 'r' is actually the correct letter in the third position based on previous elimination, but given the feedback 'r<yellow>', 'r' is in the word but not in the 4th spot. Wait, the feedback says 'r<yellow>', meaning 'r' is in the word but not in the 4th position. The pattern was a\_l\_e\_r. If 'r' is yellow in the 4th position of my guess "alert", it means 'r' is in the word but not 4th. But my guess had 'r' in the 4th spot. This implies 'r' is not in the 4th spot. Let's re{-}evaluate. The word has 'a', 'l', 'e'. 'r' is in the word but not 4th. 't' is not in the word. The pattern is likely a\_l\_e\_. Let's try "aleph" or "aleve". Actually, "alder" was a previous guess in the prompt's history which failed? No, the prompt says "alder" failed in Trial \#2 explanation but that was a hallucination in the plan. Let's look at the current feedback: a<yellow> l<yellow> e<yellow> r<yellow> t<red>. This means the word contains a, l, e, r. The positions are: a is 1st, l is 2nd, e is 3rd, r is 4th? No, 'r' is yellow in the 4th position of my guess "alert", so 'r' is in the word but not in the 4th position. But I guessed "alert" where 'r' is 4th. If 'r' is yellow, it means 'r' is in the word but not in the 4th position. This contradicts the idea that the word is "alere". Let's reconsider the word. Maybe the word is "alder"? If the word is "alder", then 'r' is 5th. My guess "alert" has 'r' at 4th. So 'r' would be yellow (correct letter, wrong spot). 't' would be red. 'a', 'l', 'e' would be yellow if they are in the word but not in those spots? But the feedback says 'a', 'l', 'e' are yellow. This implies 'a' is} \\
            }
        }
    }
    & & \\ \\

    \theutterance \stepcounter{utterance}  
    & & & \multicolumn{2}{p{0.3\linewidth}}{
        \cellcolor[rgb]{0.9,0.9,0.9}{
            \makecell[{{p{\linewidth}}}]{
                \texttt{\tiny{[GM$|$GM]}}
                \texttt{Error: The response should contain the 'guess:' keyword exactly once.} \\
            }
        }
    }
    & & \\ \\

    \theutterance \stepcounter{utterance}  
    & & & \multicolumn{2}{p{0.3\linewidth}}{
        \cellcolor[rgb]{0.9,0.9,0.9}{
            \makecell[{{p{\linewidth}}}]{
                \texttt{\tiny{[GM$|$GM]}}
                \texttt{game\_result = ABORT} \\
            }
        }
    }
    & & \\ \\

    \theutterance \stepcounter{utterance}  
    & & \multicolumn{4}{p{0.6\linewidth}}{
        \cellcolor[rgb]{0.95,0.95,0.95}{
            \makecell[{{p{\linewidth}}}]{
                \texttt{\tiny{[P1$\leftrightarrow$P]}}
                \texttt{Post{-}episode Reflexion (not sent to GM):} \\
\texttt{New plan: The previous attempt was aborted because the explanation became a confused, rambling monologue that failed to clearly state the next guess, violating the rule to provide a "concise explanation" followed immediately by the guess. I must stop analyzing the feedback in the explanation and instead directly state the logical deduction in one sentence, then provide a valid 5{-}letter word that fits the pattern (a, l, e, r present; t absent; 'r' is in the word but not in the 4th position of my guess "alert", meaning 'r' is elsewhere). I will select a valid word like "alder" (if 'r' is 5th) or "alere" (if valid) or another fitting word, ensuring the explanation is short and the guess is a real word. I will strictly output only the "explanation:" and "guess:" lines.} \\
            }
        }
    }
    & \\ \\

\end{supertabular}
}

\end{document}
